% BEGIN REV09_Q_COMPLEJO
\subsection{Iteration of the ratio reader and phase recovery}
\label{corpus9:iv:q}

The reader \(R(x,y)=x/y^2\), applied to the positive representations
of closure and self-scaling, admits the oriented composition
\[
 \mathcal Q_{\mathrm{rat}}(x,y)
   =\bigl(R(x,y),R(y,x)\bigr)
   =(x/y^2,y/x^2),\qquad x,y>0.
\]
This composition receives coordinates already produced by
APP--TRIT--TPK. Its study retains both the product \(P=xy\)
and the oriented ratio \(O=x/y\), whose transformations are
\[
 P\circ\mathcal Q_{\mathrm{rat}}=P^{-1},\qquad
 O\circ\mathcal Q_{\mathrm{rat}}=O^3.
\]
The notation \(\mathcal Q_{\mathrm{rat}}\) distinguishes this map
from the interval extension of the scalar reader.

\begin{proposition}[Positive iteration and complex extension]
\label{corpus9:iv:q-iteracion}
For \(n\geq0\), let
\[
 a_n=\frac{(-1)^n+3^n}{2},\qquad
 b_n=\frac{(-1)^n-3^n}{2}.
\]
Then
\begin{equation}
 \mathcal Q_{\mathrm{rat}}^n(x,y)
   =(x^{a_n}y^{b_n},x^{b_n}y^{a_n}).
 \label{corpus9:iv:q-potencias}
\end{equation}
The same Laurent expression defines a map on
\((\mathbb C^\times)^2\). For \(n\geq1\), this map is an
unramified covering of degree \(3^n\), with kernel
\begin{equation}
 \ker\mathcal Q_{\mathrm{rat}}^n
   =\{(\zeta,\zeta^{-1}):\zeta^{3^n}=1\}.
 \label{corpus9:iv:q-nucleo}
\end{equation}
Its restriction to the positive quadrant has a unique positive
preimage for every positive pair.
\end{proposition}
\begin{proof}
In positive logarithmic coordinates, the matrix is
\[
 B=\begin{pmatrix}1&-2\\-2&1\end{pmatrix}.
\]
The vectors \((1,1)\) and \((1,-1)\) have eigenvalues
\(-1\) and \(3\). Decomposition along these two directions gives
\(B^n=\left(\begin{smallmatrix}a_n&b_n\\b_n&a_n\end{smallmatrix}\right)\)
and proves \eqref{corpus9:iv:q-potencias}. Since \(a_n,b_n\) are
integers, the identity extends as a Laurent identity.

Write \(m=3^n\), \(\epsilon=(-1)^n\). For an image
\((u,v)\in(\mathbb C^\times)^2\), the product of a preimage
is fixed by \(p=xy=(uv)^\epsilon\). The equations are equivalent to
\[
 x^m=u\,p^{-b_n},\qquad y=p/x.
\]
The first has exactly \(m\) distinct roots. The second
determines \(y\); moreover, the product of the two images is
\(p^\epsilon=uv\), so the second original equation also
holds. Every fibre contains exactly \(m\) points. For
\((u,v)=(1,1)\), this gives \eqref{corpus9:iv:q-nucleo}; the other
fibres are its translates in the multiplicative group.

The differential satisfies
\[
 \det D\mathcal Q_{\mathrm{rat}}^n
   =(-3)^n(xy)^{\epsilon-1}\ne0.
\]
The local description by roots of a nonzero number provides
the \(m\) local branches and proves the covering assertion.
In the positive quadrant only one of these roots is positive.
In particular, for one step,
\[
 x=(uv^2)^{-1/3},\qquad y=(u^2v)^{-1/3},
\]
with the positive real roots.
\end{proof}

In the complex extension, the condition \(xy=(uv)^{-1}\)
couples the two cube roots in the one-step inverse. The
coordinates \((P,O)\) also identify
\((x,y)\) and \((-x,-y)\); their map has kernel
\(\{(1,1),(-1,-1)\}\). The preceding proof retains \((x,y)\)
and therefore the distinction between this coordinate
identification and the fibres of \(\mathcal Q_{\mathrm{rat}}^n\).

\begin{proposition}[Reconstruction from the ternary register]
\label{corpus9:iv:q-acarreo}
On the subgroup
\[
 (x,y)=(e^{2\pi i\theta},e^{-2\pi i\theta}),
\]
the map induces \(\theta\mapsto3\theta\pmod1\).
With representatives \(\theta_j\in[0,1)\), let
\[
 d_j=\lfloor3\theta_j\rfloor,\qquad
 \theta_{j+1}=3\theta_j-d_j,\qquad
 C_n=\sum_{j=0}^{n-1}d_j3^{n-1-j}.
\]
Exact recovery is given by
\begin{equation}
 \theta_0=\frac{\theta_n+C_n}{3^n}.
 \label{corpus9:iv:q-recuperacion}
\end{equation}
\end{proposition}
\begin{proof}
Substitution into the Laurent expression gives phase tripling.
The recurrence \(C_{n+1}=3C_n+d_n\), with \(C_0=0\),
gives \(\theta_n=3^n\theta_0-C_n\) by induction. This proves
\eqref{corpus9:iv:q-recuperacion}. Each \(d_j\) belongs to
\(\{0,1,2\}\), so \(n\) ternary digits retain
the choice among the \(3^n\) preimages of the reduced phase.
The branches \(B_d(\theta)=(\theta+d)/3\) satisfy
\[
 B_a\circ B_b(\theta)=\frac{\theta+b+3a}{9}.
\]
Thus two steps retain the ordered digit \(3d_0+d_1\).
\end{proof}

The product returns after two iterations, whereas the oriented
ratio is raised to the ninth power. The recovery proved here
has this reader and its phase register as its domain. Its insertion
into a representation of the complete enriched state additionally
retains the maps of sheets, orientation, memory and
survival belonging to that representation.
% END REV09_Q_COMPLEJO
